% !TEX program = lualatex
% Transparencias de Fundamentos de las Matemáticas I
% Clase 08: Aplicaciones I

\documentclass[aspectratio=169,11pt]{beamer}

\usepackage{fontspec}
\usepackage[spanish,es-nodecimaldot]{babel}
\usepackage{amsmath,amssymb,mathtools}
\usepackage{booktabs}
\usepackage{csquotes}
\usepackage{xcolor}

\usetheme{Madrid}
\usecolortheme{default}
\setbeamertemplate{navigation symbols}{}
\setbeamertemplate{footline}[frame number]

\definecolor{FdMBlue}{RGB}{20,55,100}
\definecolor{FdMRed}{RGB}{130,30,30}
\definecolor{FdMGreen}{RGB}{25,100,60}

\setbeamercolor{structure}{fg=FdMBlue}
\setbeamercolor{title}{fg=white,bg=FdMBlue}
\setbeamercolor{frametitle}{fg=white,bg=FdMBlue}
\setbeamercolor{block title}{fg=white,bg=FdMBlue}
\setbeamercolor{block body}{bg=blue!3}
\setbeamercolor{alerted text}{fg=FdMRed}

\setmainfont{Latin Modern Roman}
\setsansfont{Latin Modern Sans}
\setmonofont{Latin Modern Mono}

\newcommand{\N}{\mathbb{N}}
\newcommand{\Z}{\mathbb{Z}}
\newcommand{\Q}{\mathbb{Q}}
\newcommand{\R}{\mathbb{R}}
\newcommand{\C}{\mathbb{C}}
\newcommand{\Pow}{\mathcal{P}}
\newcommand{\id}{\operatorname{id}}
\newcommand{\mcd}{\operatorname{mcd}}
\newcommand{\mcm}{\operatorname{mcm}}

\title[FdM I -- Clase 08]{Aplicaciones I}
\subtitle{Correspondencias, funciones, dominio e imagen}
\author{Antonio Falcó Montesinos}
\institute{Universidad CEU Cardenal Herrera}
\date{Fundamentos de las Matemáticas I}

\begin{document}

\begin{frame}
  \titlepage
\end{frame}

\begin{frame}{Objetivos de la clase}
\begin{itemize}
  \item Distinguir correspondencias y aplicaciones.
  \item Definir dominio, codominio, imagen e imagen recíproca.
  \item Leer correctamente la notación \(f:A\to B\).
\end{itemize}
\end{frame}

\begin{frame}{Mapa de la clase}
\tableofcontents
\end{frame}

\section{De correspondencias a aplicaciones}

\begin{frame}{De correspondencias a aplicaciones}
\begin{block}{Correspondencia}
\begin{itemize}
  \item Una correspondencia relaciona elementos de un conjunto con elementos de otro.
  \item Puede dejar elementos sin imagen o asignar varias imágenes.
  \item Es una noción más flexible que la de aplicación.
\end{itemize}
\end{block}
\end{frame}

\begin{frame}{De correspondencias a aplicaciones}
\begin{block}{Aplicación}
\begin{itemize}
  \item Una aplicación \(f:A\to B\) asigna a cada \(a\in A\) un único elemento de \(B\).
  \item \(A\) es el dominio y \(B\) es el codominio.
  \item La unicidad de la imagen es parte de la definición.
\end{itemize}
\end{block}
\end{frame}

\begin{frame}{De correspondencias a aplicaciones}
\begin{block}{Notación}
\begin{itemize}
  \item La expresión \(f(a)\) denota la imagen de \(a\).
  \item No debe confundirse la aplicación \(f\) con el valor \(f(a)\).
  \item El codominio forma parte de los datos de la aplicación.
\end{itemize}
\end{block}
\end{frame}

\section{Imagen e imagen recíproca}

\begin{frame}{Imagen e imagen recíproca}
\begin{block}{Imagen}
\begin{itemize}
  \item La imagen de \(f\) es \(\operatorname{Im}(f)=\{f(a):a\in A\}\).
  \item Siempre se cumple \(\operatorname{Im}(f)\subseteq B\).
  \item La imagen puede ser menor que el codominio.
\end{itemize}
\end{block}
\end{frame}

\begin{frame}{Imagen e imagen recíproca}
\begin{block}{Imagen recíproca}
\begin{itemize}
  \item Si \(C\subseteq B\), entonces \(f^{-1}(C)=\{a\in A:f(a)\in C\}\).
  \item No requiere que \(f\) tenga aplicación inversa.
  \item Es una operación sobre subconjuntos.
\end{itemize}
\end{block}
\end{frame}

\section{Profundización}

\begin{frame}{Profundización}
\begin{block}{Datos de una aplicación}
\begin{itemize}
  \item Una aplicación no queda determinada solo por una fórmula.
  \item También importan el dominio y el codominio.
  \item La misma regla puede definir aplicaciones con propiedades distintas si cambia el codominio.
\end{itemize}
\end{block}
\end{frame}

\begin{frame}{Profundización}
\begin{block}{Imagen e imagen recíproca}
\begin{itemize}
  \item La imagen de un subconjunto \(S\subseteq A\) contiene los valores \(f(s)\) con \(s\in S\).
  \item La imagen recíproca de \(T\subseteq B\) contiene los elementos de \(A\) que llegan a \(T\).
  \item La imagen recíproca está definida aunque \(f\) no sea biyectiva.
\end{itemize}
\end{block}
\end{frame}

\begin{frame}{Profundización}
\begin{block}{Ejemplo de aula}
\begin{itemize}
  \item Sea \(f:\mathbb{Z}\to\mathbb{Z}\), \(f(n)=n^2\).
  \item La imagen no contiene números negativos.
  \item La imagen recíproca de \(\{4\}\) es \(\{-2,2\}\).
\end{itemize}
\end{block}
\end{frame}

\begin{frame}{Cierre}
\begin{itemize}
  \item Una aplicación exige existencia y unicidad de imagen.
  \item Dominio y codominio son parte de la definición.
  \item La imagen recíproca no es lo mismo que la inversa de una función.
\end{itemize}
\end{frame}

\begin{frame}{Trabajo recomendado}
\begin{itemize}
  \item Releer las secciones correspondientes del manual.
  \item Identificar definiciones, símbolos nuevos y enunciados principales.
  \item Preparar dudas y ejemplos para el seminario asociado.
\end{itemize}
\end{frame}

\end{document}
